\selectlanguage{english}
\pdfbookmark[0]{Overview}{companion-overview}
\noindent\textbf{Chapter XIII, first part: building the machinery that can later be forgotten.}

These sixty manuscript pages turn from the intrinsic calculus of a
localizer to the auxiliary structures that make it calculable. Grothendieck
calls a pair $(\mathcal M,W)$ a category of models when $W$ contains the
isomorphisms and satisfies two-out-of-three. This is much less than what
``model category'' ordinarily means today. He seeks classes of cofibrations
and fibrations that provide replacements, homotopy fibres and cofibres, and
eventually the two derived Kan extensions. Their purpose is to establish
properties of the localized diagram categories
\[
 \operatorname{HO}_W(I)
   =\mathcal M^{I^{\mathrm{op}}}[W(I)^{-1}],
\]
where $W(I)$ is defined objectwise. Once these constructions and their
universal properties have been established, the auxiliary classes should
cease to be part of the intrinsic answer
\citep[manuscript pp.~12--14]{Grothendieck1990XIII}.

The chapter opens with the algebra of lifting properties. In current
notation, Grothendieck's $\Phi_*$ is $\Phi^{\square}$, his $\Psi^*$ is
${}^{\square}\Psi$, and his ``couple de Quillen'' is a weak factorization
system. It is not what is now called a Quillen adjunction. The Galois
connection explains why left classes are stable under pushouts, retracts,
and continuous transfinite compositions, and why right classes have the
dual properties. The small-object theorem then supplies factorizations
from a set of generating arrows in a locally presentable category. Our
commentary proves the construction with its cardinal bound and separates
the cellular class from its retract closure. This distinction matters
because factorization can exist before the two classes are closed under
retracts.

The intrinsic classes $\operatorname{Cof}_W$ and
$\operatorname{Fib}_W$ express preservation of weak equivalences under
cobase change and base change. They are not defined by a lifting property.
Their relation to the universally trivial arrows is established using the
factorization hypotheses, not just two-out-of-three. The base-change
argument on manuscript pages~10--11 is particularly economical: factor the
change of base, apply the fibration condition along its weak-equivalence
part, and use two-out-of-three. The commentary expands this proof and gives
a three-element-poset example showing why the factorization assumption
cannot simply be omitted.

The conjectural theorem of manuscript pages~14--26 asks whether the
intrinsic classes contain small collections of tests sufficient to build
the desired structures. Here three different forms of smallness enter:
closure under filtered colimits, accessibility of a full subcategory of the
arrow category, and generation of a lifting class by a set. They do not
mean the same thing. In particular, even an accessible left lifting class
can fail to be generated by a set in the ordinary small-object sense; the
order-embedding example is worked out in the commentary. Conversely, one
specific conjecture in these pages has a precise later answer: intersections
of two set-generated left classes in a locally presentable category are
again set-generated, by the limit theorem for combinatorial cellular
categories \citep[Theorem~3.2 and Remark~3.8]{MakkaiRosicky2014Cellular}.

Smith's recognition theorem gives a useful repair of the larger programme.
One fixes a set of generating cofibrations, assumes that its right
orthogonal consists of weak equivalences, requires the trivial cofibrations
to be stable under pushouts and transfinite composition, and imposes a
solution-set condition on $W$. Under these hypotheses one obtains an
ordinary cofibrantly generated model structure
\citep[Theorem~1.7]{Beke2000}. With the intrinsic inclusions envisaged here,
the resulting structure is proper. This is a theorem with explicit extra
hypotheses, not a proof that the manuscript's filtered-colimit assumption
alone yields every assertion in its conjectural equivalence. The audit
records the individual implications and their dependencies rather than
assigning one status to the whole discussion.

The second formulation, beginning on manuscript page~27, isolates a
different issue. A weak Quillen structure is called exact when
$W=TF\circ TC$; it is left exact when $TC=C\cap W$ and right exact when
$TF=F\cap W$. These distinctions are meaningful for classes that have not
been closed under retracts. The marginal claim that retract closure of
$C,F,W$ alone makes the two one-sided exactness conditions equivalent is
false. Our example in sets takes $C=TC$ to be the monomorphisms, $F$ the
surjections, $W$ all maps, and $TF$ the surjections whose fibres are singletons
or infinite. It satisfies the weak axioms and left exactness, but excludes
the map from a two-element set to a singleton from $TF$. The missing
retract condition is visible directly.

Exactness nevertheless permits a useful completion without any retract
closure assumption. If $j\in C\cap W$, factor $j=qi$ with $i\in TC$ and
$q\in TF$. The lifting axiom makes $j$ a retract of $i$, so $j$ has the left
lifting property against every fibration. The dual argument applies to
$F\cap W$. Replacing the two designated trivial classes by $C\cap W$ and
$F\cap W$ therefore gives a biexact Quillen structure with the same
$C,F,W$. Properness supplies the required stability under base and cobase
change. This is a completion of the designated classes; the counterexample
shows why it cannot be replaced by an assertion that they were equal
already. When all four original classes are closed, the completion makes
no change, and a \emph{closed exact} weak structure is already perfect.
We use this observation again at the end of these pages to remove an
unnecessary admissibility restriction on right Quillen functors.

The discussion of homotopy on manuscript pages~29--34 concerns the actual
morphisms in the localization. Replacing objects by objects that are both
cofibrant and fibrant is not yet the same as proving that all localized
morphisms are represented by single arrows. Exactness supplies such
representatives. We first prove a description by the congruence generated
jointly by left and right homotopies. The completion then identifies each
of the usual homotopies with equality after localization for a cofibrant
source and a fibrant target. It also resolves the doubts on manuscript
pages~32--34: homotopy-cartesian and homotopy-cocartesian squares are exact
under the weak exact axioms themselves. An endpoint-lifting argument makes
the first assertion explicit; the extra hypothesis about the diagram
category used in the source's proposed proof is unnecessary for this
conclusion. Pointedness gives the suspension and loop sequences. Stability
is the additional condition needed for a triangulated structure. The
separate suggestion about raw right homotopy depends on the path-object
convention. A graph example, adapted from \citet[Example~3.6]{KurzRosicky2005},
shows that paths supplied by the designated factorization system need not
give a transitive relation without exactness, even on fibrant-cofibrant
objects. Allowing some weak-equivalence--fibration path object gives an
equivalence relation for a fibrant target by the concatenation argument
proved in the commentary.

The passage to presheaves supplies the chapter's central technical test.
Objectwise cofibrations lead to an injective diagram structure. Its
fibrations are defined by lifting against all objectwise trivial
cofibrations, and need not be merely objectwise fibrations. The calculation
\[
 (\alpha_{s!}A)(t)=\coprod_{\operatorname{Hom}_I(t,s)} A
\]
explains the coproduct hypothesis in the evaluation lemma. The subsequent
failure of componentwise factorization to solve natural lifting problems
is exhibited by explicit examples. Modern existence results for injective
structures apply in the combinatorial setting, while the cellular-category
theorem also accounts for the set-generation of the pointwise left
classes \citep[Corollary~3.4]{MakkaiRosicky2014Cellular}.

The variances then become transparent. Restriction preserves the
objectwise cofibrations and trivial cofibrations, so its right adjoint is
right Quillen. Its left adjoint is not automatically left Quillen for these
same injective structures. Above all, an ordinary Kan extension and its
derived functor are different constructions. Grothendieck explicitly
withdraws the premature identification on manuscript pages~49--50. The
commentary supplies counterexamples and the appropriate derived statements.
Under the uniform closed exact hypotheses finally adopted on manuscript
page~57, every existing ordinary right Kan extension preserves weak
equivalences between injectively fibrant diagrams by Brown's lemma.
Consequently the final admissibility conclusion holds for every such
indexing functor, beyond the finite posets singled out in the text.

Complexes give a complementary warning about size. For an infinite
indexing category, a diagram of individually bounded-below complexes need
not have a common lower bound. It therefore need not be a bounded-below
complex in the abelian category of diagrams. A family with nonzero
cohomology in successively lower degrees makes the resulting failure of a
product in $D^+(\mathbf{Ab})$ explicit. Unbounded derived operations in
Grothendieck abelian categories are instead available through suitable
resolutions and model structures; the boundedness issue concerns the
chosen subcategory and cannot be removed by renaming it.

The later derivator constructions realize the intended independence from
presentation. Cisinski supplies homotopy right Kan extensions for general
closed model categories with small limits
\citep[Theorem~6.11]{Cisinski2003}; Renaudin compares the localized
2-category of suitable combinatorial presentations with derivators of
small presentation \citep[Theorems~3.3.2 and~3.4.4]{Renaudin2009}; and the
infinity-categorical treatment recovers the same pointwise comparison
laws from Kan extensions in the underlying infinity-category
\citep[Proposition~5.3.9 and Remark~5.3.10]{RiehlVerity2017}. These results
give precise answers in their respective settings. The final
order-topological extension also becomes precise: a noetherian sober
Alexandrov $T_0$ space is finite, so under the ordinary poset topology the
apparently broader class of shapes on manuscript page~60 adds no new cases.
For the weakest recognition assertions of the manuscript, what was
established here is kept apart from what was not located in the checked
literature. An unlocated proof is not certified as an open problem.

\medskip
\noindent\textit{Commentary and audit: GPT-6 Astra (OpenAI, ChatGPT Pro
mode). This credit concerns the present contribution only; it does not
certify any subsequent assembly or editing.}
